% TOP_PROBLEM: {"id": 3305, "title": "Highly arc transitive two ended digraphs", "category_id": 3, "category": {"id": 3, "name": "graph_theory", "display_name": "Graph Theory", "description": "Problems involving graphs, networks, and their properties.", "slug": "graph-theory", "order_index": 3, "created_at": "2026-07-31T15:26:25.671Z"}, "status": "open"}
% FIRST_POSTED: null
% ENTRY_KIND: statement_only
% Imported from: batch03.tex; UnsolvedMath ID 3305
\documentclass[11pt]{article}
\usepackage[margin=25mm]{geometry}
\usepackage{fontspec}
\setmainfont{Latin Modern Roman}
\usepackage{amsmath,amssymb,mathtools}
\usepackage{unicode-math}
\setmathfont{Latin Modern Math}
\usepackage{xurl,hyperref}
\hypersetup{colorlinks=true,urlcolor=blue}
\setcounter{secnumdepth}{0}
\setlength{\parindent}{0pt}
\setlength{\parskip}{5pt}
\setlength{\emergencystretch}{3em}
\newcommand{\sref}[2]{\hyperlink{source:#1:#2}{[S#2]}}
\newcommand{\cataloguescope}[1]{\paragraph{Recorded target.}#1}
\begin{document}
% UnsolvedMath ID 3305; source catalogue code OPG-691
\setcounter{section}{3304}
\section{Highly arc transitive two ended digraphs}\label{3305}
{\small\sffamily UnsolvedMath ID 3305 · Graph Theory}
\subsection{Definitions and mathematical statement}

\paragraph{Shared notation.}$\mathbb N=\{1,2,\ldots\}$, and $\mathbb Z,\mathbb Q,\mathbb R,\mathbb C$ denote the integers, rationals, reals and complex numbers. Any other conventions are specified locally.


An end of an undirected graph is an equivalence class of one-way infinite paths, where two paths are equivalent when their tails lie in the same component after deletion of every finite vertex set. For a digraph, ends refer to the underlying undirected graph. A digraph is highly arc-transitive if its automorphisms are transitive on every set of non-backtracking directed $s$-arcs, for all $s\geq0$. In the two-ended setting of the source, choose its automorphism-invariant partition $(X_i)_{i\in\mathbb Z}$ into finite blocks, with each arc directed from $X_i$ to $X_{i+1}$. The tile at $i$ is the bipartite digraph induced between these two consecutive blocks.
\paragraph{Statement recorded in the supplied source.}
The source conjecture asks whether every tile of a two-ended highly arc-transitive digraph is a disjoint union of complete bipartite digraphs, each with all arcs oriented consistently from $X_i$ to $X_{i+1}$. The cited primary paper constructs counterexamples to this historical assertion. \sref{3305}{2} \sref{3305}{3}
\subsection{Short English statement}
Must each finite layer-to-layer piece of a two-ended highly symmetric digraph split into complete bipartite components? Counterexamples to this historical statement are known.
\subsection{Sources}
\begin{description}
\item[\hypertarget{source:3305:1}{S1}] ULAM.AI and the UnsolvedMath Contributors, UnsolvedMath: A Curated Collection of Open Mathematics Problems, record 3305 (OPG-691). CC BY 4.0. \url{https://huggingface.co/datasets/ulamai/UnsolvedMath}.
\item[\hypertarget{source:3305:2}{S2}] Open Problem Garden, Highly arc transitive two ended digraphs, node 691, original problem page. Original question and full discussion. \url{https://www.openproblemgarden.org/op/highly_arc_transitive_two_ended_digraphs}.
\item[\hypertarget{source:3305:3}{S3}] Matt DeVos, Bojan Mohar and Robert Šámal, Highly arc-transitive digraphs -- counterexamples and structure (2011), abstract. \url{https://arxiv.org/abs/1110.2945}.
\end{description}
\label{end:3305}
\end{document}
